% GENERATED FILE. Edit canonical lessons, then run npm run generate:books.
% canonical-source-hash: fnv1a64:19c0f079baef2f41
% canonical-lessons: RU-C287-prestupleniye, RU-C287-prestupnik, RU-C287-zhertva, RU-C287-ulika, RU-C287-svidetel

\chapter{Crime and its witnesses}
\label{ch:ru-crime-and-its-witnesses}
\hlchaptermodality{287}

\begin{chapteropening}
\textbf{By the end of this chapter:} \emph{I can talk about a crime: the crime, the criminal, the victim, a clue and a witness.}

Complete the last lesson of chapter 287: I can talk about a crime: the crime, the criminal, the victim, a clue and a witness.
\end{chapteropening}

\section[\texorpdfstring{\ru{преступление} (prestupléniye)}{преступление (prestupléniye)}]{\ru{преступление} (prestupléniye) — a crime}

\label{lesson:RU-C287-prestupleniye}



\begin{quote}
Before the new one: say the Russian for a publication, then the Russian for a print run.
\end{quote}

\subsection*{You'll want to know: \ru{преступление}}
\textbf{\ru{преступление}} (\emph{prestupléniye}) — \textquotedblleft{}a crime\textquotedblright{}.

\begin{quote}
\ru{Это} \ru{было} \ru{серьёзное} \ru{преступление}.
\end{quote}

\begin{quote}
\emph{Éto býlo seryóznoye prestupléniye.}
\end{quote}

\begin{quote}
It was a serious crime.
\end{quote}

Neuter. \emph{\ru{Совершить} \ru{преступление}} — to commit a crime.

\begin{cousinweb}
\emph{\ru{преступление}} is built on \emph{\ru{преступить}} \textquotedblleft{}to step over (the law)\textquotedblright{}, from \emph{\ru{пре}-}, the Church Slavonic shape of native \emph{\ru{пере}-}, and \emph{\ru{ступить}} \textquotedblleft{}to step\textquotedblright{}, the root of \emph{\ru{ступень}}, Proto-Slavic \emph{stǫpenĭ} (EtymDB).
\end{cousinweb}

\begin{grammarlens}[title={What you've built}]
One more word for telling, explaining and describing at B1.
\end{grammarlens}

\subsection*{Your turn}
\begin{itemize}
\raggedright
  \item \emph{Say it:} \emph{prestupléniye}
  \item \emph{Say it:} \emph{prestupléniye}, once more
  \item \emph{Say it:} a sentence of your own with \emph{prestupléniye}
  \item \emph{Recall:} say the Russian for a publication, then the Russian for a print run, then say \emph{prestupléniye} again
\end{itemize}

\begin{tcolorbox}[breakable,colback=teal!4,colframe=teal!35!black,title={Before you move on}]
What does \ru{преступление} mean? (\textquotedblleft{}A crime\textquotedblright{}.) Say it once more.
\end{tcolorbox}
\section[\texorpdfstring{\ru{преступник} (prestúpnik)}{преступник (prestúpnik)}]{\ru{преступник} (prestúpnik) — a criminal}

\label{lesson:RU-C287-prestupnik}



\begin{quote}
Before the new one: say the Russian for a print run, then the Russian for a crime.
\end{quote}

\subsection*{You'll want to know: \ru{преступник}}
\textbf{\ru{преступник}} (\emph{prestúpnik}) — \textquotedblleft{}a criminal\textquotedblright{}.

\begin{quote}
\ru{Преступник} \ru{убежал} \ru{через} \ru{окно}.
\end{quote}

\begin{quote}
\emph{Prestúpnik ubezhál chérez oknó.}
\end{quote}

\begin{quote}
The criminal ran away through the window.
\end{quote}

A woman is \emph{\ru{преступница}}.

\begin{cousinweb}
\emph{\ru{преступник}} is built on \emph{\ru{преступить}} \textquotedblleft{}to step over\textquotedblright{}, Church Slavonic \emph{\ru{пре}-} with \emph{\ru{ступить}} \textquotedblleft{}to step\textquotedblright{} (EtymDB lists \emph{\ru{ступень}}, Proto-Slavic \emph{stǫpenĭ}), and the ending \emph{-\ru{ник}} for a person.
\end{cousinweb}

\begin{grammarlens}[title={What you've built}]
One more word for telling, explaining and describing at B1.
\end{grammarlens}

\subsection*{Your turn}
\begin{itemize}
\raggedright
  \item \emph{Say it:} \emph{prestúpnik}
  \item \emph{Say it:} \emph{prestúpnik}, once more
  \item \emph{Say it:} a sentence of your own with \emph{prestúpnik}
  \item \emph{Recall:} say the Russian for a print run, then the Russian for a crime, then say \emph{prestúpnik} again
\end{itemize}

\begin{tcolorbox}[breakable,colback=teal!4,colframe=teal!35!black,title={Before you move on}]
What does \ru{преступник} mean? (\textquotedblleft{}A criminal\textquotedblright{}.) Say it once more.
\end{tcolorbox}
\section[\texorpdfstring{\ru{жертва} (zhértva)}{жертва (zhértva)}]{\ru{жертва} (zhértva) — a victim; a sacrifice}

\label{lesson:RU-C287-zhertva}



\begin{quote}
Before the new one: say the Russian for a crime, then the Russian for a criminal.
\end{quote}

\subsection*{You'll want to know: \ru{жертва}}
\textbf{\ru{жертва}} (\emph{zhértva}) — \textquotedblleft{}a victim; a sacrifice\textquotedblright{}.

\begin{quote}
\ru{Жертва} \ru{ничего} \ru{не} \ru{помнила}.
\end{quote}

\begin{quote}
\emph{Zhértva nichevó ne pómnila.}
\end{quote}

\begin{quote}
The victim remembered nothing.
\end{quote}

Feminine, even of a man: \emph{\ru{он} \ru{стал} \ru{жертвой} \ru{обмана}} — he fell victim to a fraud.

\begin{cousinweb}
\emph{\ru{жертва}} is inherited from Proto-Slavic \emph{žĭrtva} \textquotedblleft{}a sacrifice\textquotedblright{}, built on \emph{žĭrti} \textquotedblleft{}to offer, to sacrifice\textquotedblright{}, which EtymDB traces to PIE \emph{g\textsuperscript{w}erH-}.
\end{cousinweb}

\begin{grammarlens}[title={What you've built}]
One more word for telling, explaining and describing at B1.
\end{grammarlens}

\subsection*{Your turn}
\begin{itemize}
\raggedright
  \item \emph{Say it:} \emph{zhértva}
  \item \emph{Say it:} \emph{zhértva}, once more
  \item \emph{Say it:} a sentence of your own with \emph{zhértva}
  \item \emph{Recall:} say the Russian for a crime, then the Russian for a criminal, then say \emph{zhértva} again
\end{itemize}

\begin{tcolorbox}[breakable,colback=teal!4,colframe=teal!35!black,title={Before you move on}]
What does \ru{жертва} mean? (\textquotedblleft{}A victim; a sacrifice\textquotedblright{}.) Say it once more.
\end{tcolorbox}
\section[\texorpdfstring{\ru{улика} (ulíka)}{улика (ulíka)}]{\ru{улика} (ulíka) — a clue, a piece of evidence}

\label{lesson:RU-C287-ulika}



\begin{quote}
Before the new one: say the Russian for a criminal, then the Russian for a victim.
\end{quote}

\subsection*{You'll want to know: \ru{улика}}
\textbf{\ru{улика}} (\emph{ulíka}) — \textquotedblleft{}a clue, a piece of evidence\textquotedblright{}.

\begin{quote}
\ru{Полицейский} \ru{нашёл} \ru{важную} \ru{улику}.
\end{quote}

\begin{quote}
\emph{Politséyskiy nashyól vázhnuyu ulíku.}
\end{quote}

\begin{quote}
The policeman found an important clue.
\end{quote}

Evidence as a whole is plural: \emph{\ru{улики}}.

\begin{cousinweb}
\emph{\ru{улика}} is built on \emph{\ru{уличить}} \textquotedblleft{}to catch out, to expose\textquotedblright{}, which is usually traced to \emph{\ru{лик}} \textquotedblleft{}face\textquotedblright{}: to bring face to face with the facts. EtymDB lists no parent.
\end{cousinweb}

\begin{grammarlens}[title={What you've built}]
One more word for telling, explaining and describing at B1.
\end{grammarlens}

\subsection*{Your turn}
\begin{itemize}
\raggedright
  \item \emph{Say it:} \emph{ulíka}
  \item \emph{Say it:} \emph{ulíka}, once more
  \item \emph{Say it:} a sentence of your own with \emph{ulíka}
  \item \emph{Recall:} say the Russian for a criminal, then the Russian for a victim, then say \emph{ulíka} again
\end{itemize}

\begin{tcolorbox}[breakable,colback=teal!4,colframe=teal!35!black,title={Before you move on}]
What does \ru{улика} mean? (\textquotedblleft{}A clue, a piece of evidence\textquotedblright{}.) Say it once more.
\end{tcolorbox}
\section[\texorpdfstring{\ru{свидетель} (svidétel)}{свидетель (svidétel)}]{\ru{свидетель} (svidétel) — a witness}

\label{lesson:RU-C287-svidetel}



\begin{quote}
Before the new one: say the Russian for a victim, then the Russian for a clue.
\end{quote}

\subsection*{You'll want to know: \ru{свидетель}}
\textbf{\ru{свидетель}} (\emph{svidétel}) — \textquotedblleft{}a witness\textquotedblright{}.

\begin{quote}
\ru{Свидетель} \ru{описал} \ru{машину}.
\end{quote}

\begin{quote}
\emph{Svidétel opisál mashínu.}
\end{quote}

\begin{quote}
The witness described the car.
\end{quote}

A woman is \emph{\ru{свидетельница}}. At a wedding, \emph{\ru{свидетель}} is also the best man.

\begin{cousinweb}
\emph{\ru{свидетель}} is inherited from Old East Slavic \emph{sŭvědětelĭ} (EtymDB), built on \emph{\ru{съ}-} and the root of \emph{\ru{ведать}} \textquotedblleft{}to know\textquotedblright{}, PIE \emph{weyd-}: one who knows what happened.
\end{cousinweb}

\begin{grammarlens}[title={What you've built}]
One more word for telling, explaining and describing at B1.
\end{grammarlens}

\subsection*{Your turn}
\begin{itemize}
\raggedright
  \item \emph{Say it:} \emph{svidétel}
  \item \emph{Say it:} \emph{svidétel}, once more
  \item \emph{Say it:} a sentence of your own with \emph{svidétel}
  \item \emph{Recall:} say the Russian for a victim, then the Russian for a clue, then say \emph{svidétel} again
\end{itemize}

\begin{tcolorbox}[breakable,colback=teal!4,colframe=teal!35!black,title={Before you move on}]
What does \ru{свидетель} mean? (\textquotedblleft{}A witness\textquotedblright{}.) Say it once more.
\end{tcolorbox}
